\appendixitem{紧群的表示}

设 \(G\) 是一个紧拓扑群，且 \(\rho : G \to \mathbf{GL}(W)\) 是 \(G\) 在有限维复向量空间上的连续表示。我们称 \(\rho\) 是实型、复型或四元数型不可约表示，当且仅当对于 \(\mathbf{C}^{(G)}\)-模 \(W\)（相对于代数 \(A = \mathbf{R}^{(G)}\)），情形也是如此。设 \(H\) 是 \(W\) 上一个分离的正埃尔米特型，且在 \(G\) 下不变。

命题 3. 假设 \(\rho\) 是不可约的。

\begin{enumerate}
\def\labelenumi{\alph{enumi})}
\item
  表示 \(\rho\) 属于实型，当且仅当 W 上存在一个非零的、在 G 下不变的对称双线性型 B。在此情形，型 B 是分离的；由满足 \(w \in W\) 且 \(\mathrm{H}(w, x) = \mathrm{B}(w, x)\) 对所有 \(x \in W\) 成立的元素构成的集合 V，是 W 上一个在 G 下不变的 \(\mathbf{R}\)-结构。
\item
  表示 \(\rho\) 属于复型，当且仅当 W 上不存在在 G 下不变的非零双线性型。
\item
  表示 \(\rho\) 属于四元数型，当且仅当 W 上存在一个在 G 下不变的非零交错双线性型；这样的型必然是分离的。
\end{enumerate}

对于 \(\theta \in \mathrm{Hom}_{\mathbf{C}^{(\mathrm{G})}}(\mathrm{W},\overline{\mathrm{W}})\) 和 \(x,y\in \mathrm{W}\)，令 \(\mathrm{B}_{\theta}(x,y) = \mathrm{H}(\theta x,y)\)。则 \(\mathrm{B}_{\theta}\) 是 \(\mathrm{W}\) 上的双线性型，在 \(\mathrm{G}\) 下不变，并且当 \(\theta\) 非零时是分离的。记 \(\mathcal{B}(\mathrm{W})^{\mathrm{G}}\) 为 \(\mathrm{W}\) 上在 \(\mathrm{G}\) 下不变的双线性型空间；映射 \(\theta \mapsto \mathrm{B}_{\theta}\) 从 \(\mathrm{Hom}_{\mathbf{C}^{(\mathrm{G})}}(\mathrm{W},\overline{\mathrm{W}})\) 到 \(\mathcal{B}(\mathrm{W})^{\mathrm{G}}\) 是 \(\mathbf{C}\)-向量空间的一个同构。特别地，这推出断言 \(b\)）。

令 \(\theta\) 是一个 \(\mathbf{C}^{(\mathrm{G})}\)-同构，从 W 到 \(\overline{\mathrm{W}}\)，满足 \(\theta \circ \theta = \alpha_{\mathrm{W}}\)，其中 \(\alpha \in \{-1, +1\}\)（命题 2）；由于 \(\mathcal{B}(\mathrm{W})^{\mathrm{G}}\) 的维数为 1，存在 \(\varepsilon \in \mathbf{C}\)，使得

\[
\mathrm{B} _ {\theta} (y, x) = \varepsilon \mathrm{B} _ {\theta} (x, y) \quad \text { for   all } x, y \text { in } \mathrm{W}.
\]

反复交换变量，得到 \(\mathrm{B}_{\theta}(y,x)=\varepsilon\mathrm{B}_{\theta}(x,y)=\varepsilon^{2}\mathrm{B}_{\theta}(y,x)\)，因此 \(\varepsilon^{2}=1\)，且 \(\varepsilon\in\{-1,+1\}\)。此外，对于 W 中的 x，

\[
\mathrm{H} (\theta x, \theta x) = \mathrm{B} _ {\theta} (x, \theta x) = \varepsilon \mathrm{B} _ {\theta} (\theta x, x) = \varepsilon \mathrm{H} (\theta \circ \theta (x), x) = \varepsilon \alpha \mathrm{H} (x, x)
\]

所以由于 H 是正的，有 \(\varepsilon\alpha > 0\)，亦即 \(\varepsilon = \alpha\)。于是断言 a) 和 c) 由命题 2 得出。

记 G 上总质量为 1 的 Haar 测度为 dg。

引理 1. 设 \(W^{G}\) 是 W 的子空间，由在 G 下不变的元素组成。W 的自同态 \(\int_{G}\rho(g)dg\) 是一个投影，其像为 \(W^{G}\)，并且与 G 的作用相容。特别地，

\[
\dim \mathrm{W} ^ {\mathrm{G}} = \int_ {\mathrm{G}} \operatorname{Tr} \rho (g) d g.
\]

令 \(p = \int_{\mathrm{G}} \rho(g) dg\)；对于 \(h \in \mathrm{G}\)，

\[
\rho (h) \circ p = \int_ {\mathrm{G}} \rho (h g) d g = \int_ {\mathrm{G}} \rho (g) d g = p
\]

类似地，\(p \circ \rho(h) = p\)。因此，p 与 G 的作用相容，且其像包含于 \(W^{G}\)。若 \(w \in W^{G}\)，则 \(p(w) = \int_{\mathbf{G}} \rho(g) w \, dg = w\)，引理得证。

引理 2. 设 \(u\) 是有限维向量空间 \(E\)（定义在域 \(K\) 上）的一个自同态。则

\[
\operatorname{Tr} u ^ {2} = \operatorname{Tr} \mathbf {S} ^ {2} (u) - \operatorname{Tr} \Lambda^ {2} (u).
\]

设 \(\chi_u(\mathrm{X}) = \prod_{i=1}^{n} (\mathrm{X} - \alpha_i)\) 是在 K 的某个适当扩域中将 \(u\) 的特征多项式分解为线性因子的一个分解。我们有 \(\operatorname{Tr} u^2 = \sum_i \alpha_i^2\)、\(\operatorname{Tr} \Lambda^2(u) = \sum_{i < j} \alpha_i \alpha_j\)、\(\operatorname{Tr} \mathbf{S}^2(u) = \sum_{i \leq j} \alpha_i \alpha_j\)（参见~《代数》，第 VII 章，§5，节 5，推论 3），结果得证。

命题 4. 假设 \(\rho\) 是不可约的。那么，\(\rho\) 属于实型（分别为复型、四元数型），当且仅当积分 \(\int_{\mathrm{G}} \mathrm{Tr} \rho(g^2) dg\) 等于 1（分别为 0、-1）。

记 \(\check{\rho}\) 为 \(\rho\) 在 \(W^{*}\) 上的反步表示（定义为 \(\check{\rho}(g)={}^{t}\rho(g^{-1})\)）。将引理 2 应用于 \(\check{\rho}(g)\)，并在 G 上积分，得到

\[
\int_ {\mathrm{G}} \operatorname{Tr} \rho (g ^ {2}) d g = \int_ {\mathrm{G}} \operatorname{Tr} ^ {t} \rho (g ^ {- 2}) d g = \int_ {\mathrm{G}} \operatorname{Tr} \mathbf {S} ^ {2} (\check {\rho} (g)) d g - \int_ {\mathrm{G}} \operatorname{Tr} \Lambda^ {2} (\check {\rho} (g)) d g
\]

故由引理 1，

\[
\int_ {\mathrm{G}} \mathrm{Tr} \rho (g ^ {2}) d g = \dim (\mathbf {S} ^ {2} \mathrm{W} ^ {*}) ^ {\mathrm{G}} - \dim (\Lambda^ {2} \mathrm{W} ^ {*}) ^ {\mathrm{G}}.
\]

但是，\(S^{2}W^{*}\)（分别为 \(\Lambda^{2}W^{*}\)）可与 W 上的对称（分别为交错）双线性型空间相识别。因此，命题立即由命题 3 得出。

\specialchapter{习题}

\exercisesection{\S~1}

\begin{enumerate}
\def\labelenumi{\arabic{enumi})}
\tightlist
\item
  设 \(\mathrm{G}\) 是一个有限维、连通、交换的复李群，且 \(\mathrm{V}\) 是其李代数。
\end{enumerate}

\begin{enumerate}
\def\labelenumi{\alph{enumi})}
\item
  映射 \(\exp_{\mathrm{G}}: \mathrm{V} \to \mathrm{G}\) 是满同态；其核 \(\Gamma\) 是 \(\mathrm{V}\) 的离散子群。
\item
  G 紧当且仅当 \(\Gamma\) 是 V 中的一个格；此时称 G 为复环面。
\item
  设 \(\Gamma\) 是 \(\mathbf{C}^2\) 的离散子群，由元素 \(e_1 = (1,0), e_2 = (0,1), e_3 = (\sqrt{2},i)\) 生成；令 \(\mathrm{G} = \mathrm{C}^2 / \Gamma\)，\(\mathrm{H} = (\Gamma + \mathrm{C}e_1) / \Gamma\)。证明 H 同构于 \(\mathbf{C}^*\)，G/H 是维数为 1 的复环面，但 G 不包含非零复环面。
\item
  每个有限维连通紧复李群都是复环面（参见~第 III 章，§6，节 3，命题 6）。
\end{enumerate}

\begin{enumerate}
\def\labelenumi{\arabic{enumi})}
\setcounter{enumi}{1}
\tightlist
\item
  设 H 是由复数 \(\tau\) 中满足 \(\mathcal{I}(\tau) > 0\) 者组成的集合。
\end{enumerate}

\begin{enumerate}
\def\labelenumi{\alph{enumi})}
\item
  证明可以在 H 上定义离散群 \(\mathbf{SL}(2,\mathbf{Z})\) 的解析左作用律，规定 \(\gamma\tau=\frac{a\tau+b}{c\tau+d}\)；对所有 \(\gamma=\begin{pmatrix}a&b\\ c&d\end{pmatrix}\in\mathbf{SL}(2,\mathbf{Z})\) 和 \(\tau\in H\)，该规定成立。
\item
  对于 \(\tau\in H\)，记 \(T_{\tau}\) 为复环面 \(\mathbf{C}/(\mathbf{Z}+\mathbf{Z}\tau)\)。证明映射 \(\tau\mapsto T_{\tau}\) 通过取商诱导出从集合 \(\mathrm{H}/\mathrm{SL}(2,\mathbf{Z})\) 到维数为 1 的复环面的同构类集合的一个双射。
\end{enumerate}

\begin{enumerate}
\def\labelenumi{\arabic{enumi})}
\setcounter{enumi}{2}
\item
  设 G 是 \(\mathbf{O}(n)\) 的一个整子群，使得 G 在 \(\mathbf{R}^{n}\) 上的恒等表示是不可约的。证明 G 是闭的（将 G 写成 \(K \times N\) 的形式，其中 K 是紧的、N 是交换的，并证明 \(\dim \overline{N} \leq 1\)）。
\item
  证明命题 3（节 3）中的条件分别等价于下列每个条件：
\end{enumerate}

(ii') 群 \(\operatorname{Ad}(\mathbf{G})\) 在 \(\operatorname{Aut}(\mathbf{L}(\mathbf{G}))\) 中相对紧。

(ii'\,') 群 \(\operatorname{Ad}(\mathbf{G})\) 在 \(\operatorname{End}(\mathbf{L}(\mathbf{G}))\) 中相对紧。

\begin{enumerate}
\def\labelenumi{(\alph{enumi})}
\setcounter{enumi}{21}
\tightlist
\item
  G 的单位元 e 的每个邻域都包含一个在内自同构下稳定的 e 的邻域（参见~《积分》，§3，节 1，命题 1）。
\end{enumerate}

\begin{enumerate}
\def\labelenumi{\arabic{enumi})}
\setcounter{enumi}{4}
\tightlist
\item
  设 A 是 \(\mathbf{GL}(3,\mathbf{R})\) 的闭子群，由矩阵 \((a_{ij})\) 组成，这些矩阵满足 \(a_{ij}=0\)（当 i\textgreater j 时）、\(a_{ii}=1\)（\(1\leq i\leq3\)）、\(a_{12}\in \mathbf{Z}\)、\(a_{23}\in \mathbf{Z}\)；令 \(\theta\in \mathbf{R}-\mathbf{Q}\)，并令 B 是 A 的子群，由矩阵 \((a_{ij})\) 组成，这些矩阵满足 \(a_{12}=a_{23}=0\) 且 \(a_{13}\in\theta\mathbf{Z}\)，并令 G=A/B。
\end{enumerate}

\begin{enumerate}
\def\labelenumi{\alph{enumi})}
\item
  证明 G 是一个维数为 1 的李群，且其李代数是紧的。
\item
  证明 \(\mathrm{C}(\mathrm{G}) = \overline{\mathrm{D}(\mathrm{G})} = \mathrm{G}_{0}\)，并且 \(G_{0}\) 是 G 的最大紧子群。
\item
  证明 G 不是 \(G_{0}\) 与任何子群的半直积。
\end{enumerate}

\begin{enumerate}
\def\labelenumi{\arabic{enumi})}
\setcounter{enumi}{5}
\item
  证明一个（实）李代数是紧的，当且仅当它存在一组基 \((e_{\lambda})_{\lambda\in\mathrm{L}}\)，使得结构常数 \(\gamma_{\lambda\mu\nu}\) 关于 \(\lambda,\mu,\nu\) 反对称（即它们满足 \(\gamma_{\lambda\mu\nu}=-\gamma_{\mu\lambda\nu}=-\gamma_{\lambda\nu\mu}\)）。
\item
  对合李代数是指一个配备自同构 s 的（实）李代数 \(\mathfrak{g}\)，满足 \(s \circ s = 1_{\mathfrak{g}}\)。记 \(\mathfrak{g}^{+}\)（分别为 \(\mathfrak{g}^{-}\)）为 \(\mathfrak{g}\) 中 s 对应于特征值 +1（分别为 -1）的特征子空间。
\end{enumerate}

\begin{enumerate}
\def\labelenumi{\alph{enumi})}
\item
  证明 \(\mathfrak{g}^{+}\) 是 \(\mathfrak{g}\) 的子代数，且 \(\mathfrak{g}^{-}\) 是一个 \(\mathfrak{g}^{+}\)-模；有 \([\mathfrak{g}^{-},\mathfrak{g}^{-}] \subset \mathfrak{g}^{+}\)，并且 \(\mathfrak{g}^{+}\) 与 \(\mathfrak{g}^{-}\) 关于 Killing 型正交。
\item
  证明下列条件等价：
\end{enumerate}

\begin{enumerate}
\def\labelenumi{(\roman{enumi})}
\item
  \(\mathfrak{g}^{+}\)-模 \(\mathfrak{g}^{-}\) 是单模。
\item
  代数 \(\mathfrak{g}^{+}\) 在所有不同于 \(\mathfrak{g}\) 的 \(\mathfrak{g}\) 子代数中是极大的。若这些条件成立，且 \(\mathfrak{g}^{+}\) 不包含 \(\mathfrak{g}\) 的非零理想，则称对合李代数 \((\mathfrak{g}, s)\) 为不可约的。
\end{enumerate}

\begin{enumerate}
\def\labelenumi{\alph{enumi})}
\setcounter{enumi}{2}
\tightlist
\item
  假设 \(\mathfrak{g}\) 是半单的。证明下列条件等价：
\end{enumerate}

\begin{enumerate}
\def\labelenumi{(\roman{enumi})}
\item
  \(\mathfrak{g}\) 的在 \(s\) 下稳定的理想只有 \(\{0\}\) 和 \(\mathfrak{g}\)；
\item
  \(\mathfrak{g}\) 或者是单的，或者是两个在 \(s\) 下互换的单理想之和。
\end{enumerate}

证明当 \((\mathfrak{g}, s)\) 不可约时这些条件成立（注意，\(\mathfrak{g}\) 是满足 (ii) 且在 s 下稳定的理想的直和）。

\begin{enumerate}
\def\labelenumi{\alph{enumi})}
\setcounter{enumi}{3}
\item
  从现在起假设 \(\mathfrak{g}\) 是紧半单的。证明对合李代数 \((\mathfrak{g}, s)\) 不可约，当且仅当 s 不等于恒等映射，且 \((\mathfrak{g}, s)\) 满足 c) 中的等价条件（令 p 为 \(\mathfrak{g}^{+}\)-模 \(\mathfrak{g}^{-}\) 的一个子模，q 为其关于 Killing 型的正交补；注意到 \([p, q] = 0\)，并由此推出 \(p + [p, p]\) 是 \(\mathfrak{g}\) 的理想）。
\item
  证明 \(\mathfrak{g}\) 是一族在 s 下稳定的理想 \((\mathfrak{g}_{i})_{0\leq i\leq n}\) 的直和，其中 \(\mathfrak{g}_{0}\) 由 s 固定，且对合代数 \((\mathfrak{g}_{i},s|\mathfrak{g}_{i})\) 在 \(1\leq i\leq n\) 时不可约。
\end{enumerate}

\begin{enumerate}
\def\labelenumi{\arabic{enumi})}
\setcounter{enumi}{7}
\tightlist
\item
  设 G 是一个紧半单李群，u 是 G 的一个阶为 2 的自同构，K 是 u 的不动点集的单位元连通分支，X 是流形 G/K（齐性空间 X 因而称为对称空间）。
\end{enumerate}

\begin{enumerate}
\def\labelenumi{\alph{enumi})}
\item
  证明，如果 G 是概单的，则 K 在所有不同于 G 的连通闭子群中是极大的；换言之（第 III 章，§3，习题 8），G 在 X 上的作用是本原的。对称空间 X 因而称为不可约的。
\item
  假设流形 X 是单连通的；证明它同构于若干流形的乘积，其中每个流形都同构于一个李群或一个不可约对称空间。
\end{enumerate}

\begin{enumerate}
\def\labelenumi{\arabic{enumi})}
\setcounter{enumi}{8}
\tightlist
\item
  设 \(\mathfrak{a}\) 是一个实或复李代数，且 \(G\) 是 \(\operatorname{Aut}(\mathfrak{a})\) 的一个紧子群。证明 \(\mathfrak{a}\) 有一个在 \(G\) 下稳定的 Levi 子代数（第 I 章，§6，节 8，定义 7）（归结到 \(\mathfrak{a}\) 的根基为交换的情形，并使用《积分》，第 VII 章，§3，节 2，引理 2）。
\end{enumerate}

\exercisesection{\S~2}

\begin{enumerate}
\def\labelenumi{\arabic{enumi})}
\tightlist
\item
  设 \(\mathbf{G}\) 是一个连通紧李群，\(g\) 是 \(\mathbf{G}\) 的一个元素。
\end{enumerate}

\begin{enumerate}
\def\labelenumi{\alph{enumi})}
\item
  证明存在一个整数 \(n \geq 1\)，使得中心化子 \(Z(g^{n})\) 是连通的（证明对于适当的 n，由 \(g^{n}\) 生成的闭子群是一个环面）。
\item
  假设 \(\operatorname{Ker}(\operatorname{Ad}g^n - 1)\) 的维数与 \(n\) 无关（\(n \geq 1\)）。证明 \(Z(g)\) 是连通的。
\item
  如果 \(g^n\) 对于所有 \(n \geq 1\) 都是正则元，则 \(Z(g)\) 是连通的。
\end{enumerate}

\begin{enumerate}
\def\labelenumi{\arabic{enumi})}
\setcounter{enumi}{1}
\item
  证明每个连通紧李群 G 都是其导群与一个环面的半直积（注意，若 T 是 G 的极大环面，则 \(D(G) \cap T\) 是一个环面）。
\item
  设 G 是一个紧李群，\(\mathfrak{g}\) 是其李代数，\(\mathfrak{s}\) 是 \(\mathfrak{g}\) 的向量子空间，满足对于 \(x, y, z\)（它们属于 \(\mathfrak{s}\)），有 \([x, [y, z]] \in \mathfrak{s}\)。令 \(\mathcal{L}\) 为 \(\mathfrak{g}\) 中包含于 \(\mathfrak{s}\) 的交换子代数集合。证明 G 中 \(\mathfrak{s}\) 的稳定化子的单位元连通分支在 \(\mathcal{L}\) 的极大元素集合上传递地作用（论证如定理 1 的证明）。
\item
  设 G 是一个连通紧李群，T 是 G 的极大环面，S 是 T 的子环面。分别记 \(\Sigma\)（分别为 F）为 S 在 \(W_{G}(T)\) 中的稳定化子（分别为固定子）。
\end{enumerate}

\begin{enumerate}
\def\labelenumi{\alph{enumi})}
\item
  证明群 \(\mathrm{N_G(S) / Z_G(S)}\) 同构于商 \(\Sigma /\mathrm{F}\)。
\item
  设 H 是 G 的连通闭子群，且 S 是 H 的极大环面。证明 \(W_{\mathrm{H}}(\mathrm{S})\) 的每个元素，视为 S 的自同构，都是 \(\Sigma\) 中某个元素在 S 上的限制。
\end{enumerate}

\begin{enumerate}
\def\labelenumi{\arabic{enumi})}
\setcounter{enumi}{4}
\tightlist
\item
  设 \(G\) 是一个连通紧李群，且 \(S\) 是 \(G\) 的一个环面。证明下列条件等价：
\end{enumerate}

\begin{enumerate}
\def\labelenumi{(\roman{enumi})}
\item
  S 包含于唯一的极大环面中；
\item
  \(Z_{G}(S)\) 是一个极大环面；
\item
  S 包含一个正则元。
\end{enumerate}

\begin{enumerate}
\def\labelenumi{\arabic{enumi})}
\setcounter{enumi}{5}
\tightlist
\item
  设 G 是一个连通紧李群，T 是 G 的极大环面，\(\mathfrak{g}\)（分别为 \(\mathfrak{t}\)）是 G（分别为 T）的李代数，且 \(i : \mathfrak{t} \to \mathfrak{g}\) 是典则嵌入。证明映射 \(^{t}i : \mathfrak{g}^{*} \to \mathfrak{t}^{*}\) 通过取商诱导出从 \(\mathfrak{g}^{*}/G\) 到 \(\mathfrak{t}^{*}/W_{G}(T)\) 的一个同胚。
\end{enumerate}

¶ 7) 设 X 和 Y 是两个类 \(C^{r}\) 的实流形（\(1 \leq r \leq \omega\)），它们是分离的、连通的且维数为 n；设 \(f : X \to Y\) 是一个类 \(C^{r}\) 的真态射，并配备定向 F（《微分流形与解析流形》，结果 10.2.5）。

\begin{enumerate}
\def\labelenumi{\alph{enumi})}
\item
  证明存在唯一实数 d，使得 \(\int_{X} f^{*}\alpha = d\int_{Y}\alpha\) 对 Y 上每个扭曲微分形式 \(\alpha\)（次数为 n、类 \(C^{1}\) 且具有紧支集）成立（使用《微分流形与解析流形》，结果 11.2.4）。
\item
  设 \(y \in \mathbf{Y}\) 满足 \(f\) 在 \(f^{-1}(y)\) 的每一点处都无分歧。对于 \(x \in f^{-1}(y)\)，令 \(\nu_x(f) = 1\)（分别令 \(\nu_x(f) = -1\)），若映射 \(\mathrm{F}_x\) 和 \(\tilde{f}_x\)（《微分流形与解析流形》，结果 10.2.5，例 b）从 \(\mathrm{Or}(\mathrm{T}_x(\mathrm{X}))\) 到 \(\mathrm{Or}(\mathrm{T}_y(\mathrm{Y}))\) 相同（分别为相反映射）。证明 \(d = \sum_{x \in f^{-1}(y)} \nu_x(f)\)，特别地，\(d \in \mathbf{Z}\)。
\end{enumerate}

称 d 为 f 的次数，并记为 \(\deg f\)。若 X = Y 且 X 可定向，则可以取 F 为保持 X 的定向的定向。

\begin{enumerate}
\def\labelenumi{\alph{enumi})}
\setcounter{enumi}{2}
\item
  证明，如果 \(\deg f \neq 0\)，则 \(f\) 是满射。
\item
  如果存在 \(x \in X\)，使得 \(f^{-1}(f(x)) = \{x\}\)，且 f 在 x 处无分歧，则 f 是满射。
\end{enumerate}

\begin{enumerate}
\def\labelenumi{\arabic{enumi})}
\setcounter{enumi}{7}
\tightlist
\item
  设 G 是一个连通紧李群，dg 是 G 上的归一化 Haar 测度。
\end{enumerate}

\begin{enumerate}
\def\labelenumi{\alph{enumi})}
\tightlist
\item
  令 \(f: \mathrm{G} \to \mathrm{G}\) 为一个流形态射；对于 \(g \in \mathrm{G}\)，微分 \(\mathrm{T}_g(f)\) 可通过左平移识别为从 \(\mathrm{L}(\mathrm{G})\) 到 \(\mathrm{L}(\mathrm{G})\) 的线性映射。证明下列公式
\end{enumerate}

\[
\deg f. \int_ {\mathrm{G}} \varphi (g) d g = \int_ {\mathrm{G}} \varphi (f (g)) \det \mathrm{T} _ {g} (f) d g
\]

对 G 上每个可积的（复值）函数 \(\varphi\)，以及

\[
\deg f = \int_ {\mathrm{G}} \det \mathrm{T} _ {g} (f) d g.
\]

\begin{enumerate}
\def\labelenumi{\alph{enumi})}
\setcounter{enumi}{1}
\tightlist
\item
  令 \(\psi_k: \mathbf{G} \to \mathbf{G}\) 为映射 \(g \mapsto g^k\)。证明公式
\end{enumerate}

\[
\deg \psi_ {k} = \int_ {\mathrm{G}} \det (1 + \operatorname{Ad} g + \dots + (\operatorname{Ad} g) ^ {k - 1}) d g.
\]

\begin{enumerate}
\def\labelenumi{\alph{enumi})}
\setcounter{enumi}{2}
\tightlist
\item
  令 \(\mathrm{T}\) 是 \(\mathrm{G}\) 的极大环面，且 \(\mathrm{T}_r\) 是 \(\mathrm{T}\) 的正则元集合。证明 \(\exp_{\mathrm{G}}^{-1}(\mathrm{T}_r) \subset \mathrm{L}(\mathrm{T})\)，并且 \(\psi_k^{-1}(\mathrm{T}_r) \subset \mathrm{T}_r\) 对所有 \(k \geq 1\) 成立。
\end{enumerate}

\(d\)) 证明 \(\deg \psi_{k} = k^{\dim (\mathrm{T})}\)；推出等式

\[
\int_ {\mathrm{G}} \det (1 + \operatorname{Ad} g + \dots + (\operatorname{Ad} g) ^ {k - 1}) d g = k ^ {\dim (\mathrm{T})}.
\]

\begin{enumerate}
\def\labelenumi{\arabic{enumi})}
\setcounter{enumi}{8}
\tightlist
\item
  本习题旨在给出定理 2 的另一证明。设 G 是一个连通紧李群，T 是 G 的极大环面，N 是 T 的正规化子。记 \(G \times^{N} T\) 为 \(G \times T\) 按 N 的作用 \((g, t).n = (gn, n^{-1}tn)\) 得到的商流形（《微分流形与解析流形》，结果 6.5.1）。
\end{enumerate}

\begin{enumerate}
\def\labelenumi{\alph{enumi})}
\item
  证明态射 \((g, t) \mapsto gtg^{-1}\) 从 \(G \times T\) 到 G，通过取商定义出一个解析态射 \(f : G \times^{N} T \to G\)。
\item
  设 \(\theta\) 是 T 中一个元素，其幂集在 T 中稠密（《一般拓扑》，第 VII 章，§1，节 3，命题 7 的推论 2）。证明 \(f^{-1}(\theta) = \{x\}\)，其中 \(x\) 是 \((e,\theta)\) 在 \(\mathrm{G} \times^{\mathrm{N}}\mathrm{T}\) 中的类，并且 \(f\) 在 \(x\) 处无分歧。
\item
  由习题 7 d) 得出 \(f\) 是满射，并由此给出定理 2 的另一证明。
\end{enumerate}

\begin{enumerate}
\def\labelenumi{\arabic{enumi})}
\setcounter{enumi}{9}
\tightlist
\item
  * 本习题使用代数拓扑中的下列结果（Lefschetz 公式 \(^{11}\)）：设 X 是有限维紧流形，f 是从 X 到自身的一个态射。假设 X 中满足 \(f(x) = x\) 的点 x 所成的集合 F 是有限的，并且对于所有 \(x \in F\)，数 \(\delta(x) = \det(1 - \mathrm{T}_{x}(f))\) 都非零。如果 \(\mathrm{H}^{i}(f)\) 表示由 f 诱导的向量空间 \(\mathrm{H}^{i}(\mathrm{X}, \mathbf{R})\) 的自同态（对于 \(i \geq 0\)），则
\end{enumerate}

\[
\sum_ {x \in \mathrm{F}} \delta (x) / | \delta (x) | = \sum_ {i \geq 0} (- 1) ^ {i} \operatorname{Tr} \mathrm{H} ^ {i} (f).
\]

设 G 是一个连通紧李群，T 是 G 的极大环面。对于 \(g \in G\)，记 \(\tau(g)\) 为由 g 的左乘法在流形 G/T 上诱导的自同构。

\begin{enumerate}
\def\labelenumi{\alph{enumi})}
\item
  设 \(t\) 是 \(\mathrm{T}\) 中的一个元素，使得由 \(t\) 生成的子群在 \(\mathrm{T}\) 中稠密。证明 \(\tau(t)\) 的不动点正是 \(n\mathrm{T}\)（其中 \(n \in \mathrm{N}_{\mathrm{G}}(\mathrm{T})\)）的类。从而 \((\mathrm{N}_{\mathrm{G}}(\mathrm{T}):\mathrm{T}) = \sum_{i}(-1)^{i}\dim_{\mathbf{R}}\mathrm{H}^{i}(\mathrm{G / T},\mathbf{R})\)。
\item
  令 g 是 G 的任意元素；证明 \(\tau(g)\) 的不动点集非空，并由此给出定理 2 的另一证明。*
\end{enumerate}

\begin{enumerate}
\def\labelenumi{\arabic{enumi})}
\setcounter{enumi}{10}
\tightlist
\item
  设 G 是一个紧李群。如果 G 的子群 S 等于某个在其正规化子中具有有限指数的循环子群的闭包，则称 S 为 (C) 型子群。
\end{enumerate}

\begin{enumerate}
\def\labelenumi{\alph{enumi})}
\item
  设 \(S\) 是一个 (C) 型子群；证明 \(S_0\) 是 \(G_0\) 的极大环面，且 \(S\) 是 \(S_0\) 与某个有限循环子群的直积。
\item
  证明 G 的每个元素 g 都包含于一个 (C) 型子群中（考虑由 g 和 \(Z(g)_{0}\) 的一个极大环面生成的群）。
\item
  设 S 是一个 (C) 型子群，且 s 是 S 中幂集在 S 中稠密的元素。证明 \(sG_{0}\) 中的每个元素都在 \(\operatorname{Int}(G_{0})\) 下与 \(sS_{0}\) 中的某个元素共轭（使用习题 10 的方法）。
\item
  记 \(p: G \to G/G_{0}\) 为商映射。证明映射 \(S \mapsto p(S)\) 从 G 的 (C) 型子群的共轭类集合到 \(G/G_{0}\) 的循环子群的共轭类集合诱导出一个双射。
\item
  设 S 是 G 的一个 (C) 型子群；记 \(S_{\rho}\) 为 S 中其在 \(S/S_{0}\) 中的类生成 \(S/S_{0}\) 的元素集合。证明 \(S_{\rho}\) 中在 G 中共轭的两个元素在 \(N_{G}(S)\) 中共轭。
\end{enumerate}